\documentclass[11pt]{article}
\usepackage[margin=1in]{geometry}
\usepackage{amsmath,amssymb,amsthm}
\usepackage{graphicx}
\usepackage[hidelinks]{hyperref}
\usepackage{microtype}

\newtheorem{theorem}{Theorem}
\newtheorem{remark}{Remark}

\title{The Modulus Map Refutes the Direct Non-Monotone\\
Sublinear Wulbert Analogue in $L^1[0,1]$}
\author{Run fa\_banach\_001, agent\_lane\_16}
\date{2026-08-11}

\begin{document}
\maketitle

\section*{Status}

\textbf{Candidate full counterexample to the direct analogue; likely valid,
wording-scope review needed.}  The source asks informally for ``an analogue.''
The result below completely refutes the direct formulation obtained by
retaining Wulbert's hypotheses and replacing linear maps by continuous
sublinear maps.  A modified theorem with sign-symmetric tests is not claimed
to be false.

\section*{Source question and precise direct reading}

Gal--Niculescu ask on PDF page 12 of arXiv:2206.14102 whether Wulbert's
non-monotone Korovkin theorem for operators on $L^1[0,1]$ admits an analogue
for continuous sublinear maps.

\begin{figure}[h]
\centering
\includegraphics[width=0.96\textwidth]{figures/open_question_crop.png}
\caption{The complete source question, arXiv:2206.14102v1, PDF page 12.}
\end{figure}

Wulbert's Theorem 9 is restated in Theorem 1.3 of arXiv:2403.03476.  For
linear $L_n:L^1[0,1]\to L^1[0,1]$, its hypotheses are
\begin{enumerate}
  \item $L_n(1)\to1$ in norm;
  \item $L_n(x)\to x$ and $L_n(x^2)\to x^2$ weakly;
  \item $\|L_n\|\to1$;
\end{enumerate}
and the conclusion is $\|L_nf-f\|_1\to0$ for every $f\in L^1[0,1]$.
The direct sublinear analogue changes only the word ``linear.''

\section*{Counterexample}

\begin{theorem}[Positive-test no-go]
Let $\mathcal K\subset L^1[0,1]_+$ be any family.  There is a constant
sequence of continuous sublinear operators $T_n:L^1[0,1]\to L^1[0,1]$ such
that $\|T_n\|=1$, $T_nh=h$ for every $h\in\mathcal K$, but $T_n$ does not
converge to the identity even weakly on all of $L^1[0,1]$.
\end{theorem}

\begin{proof}
For every $n$, set
\[
                 T_nf=|f|.
\]
The pointwise triangle inequality and positive homogeneity give
\[
 T_n(f+g)\le T_nf+T_ng,
 \qquad T_n(\alpha f)=\alpha T_nf\quad(\alpha\ge0),
\]
which is exactly the source's definition of sublinearity.  Also
\[
 \|T_nf-T_ng\|_1
 =\big\||f|-|g|\big\|_1\le\|f-g\|_1,
 \qquad \|T_nf\|_1=\|f\|_1.
\]
Thus $T_n$ is continuous (indeed 1-Lipschitz) and its nonlinear operator norm
is one.  Every $h\ge0$ is fixed exactly.

On the other hand,
\[
 T_n(-1)=1,
 \qquad \|T_n(-1)-(-1)\|_1=2.
\]
Pairing the difference with the constant member $1$ of $L^\infty[0,1]$ also
gives $2$, so even weak convergence fails.
\end{proof}

Taking $\mathcal K=\{1,x,x^2\}$ verifies all three Wulbert hypotheses in
strictly stronger form: each test is fixed in norm and $\|T_n\|=1$ exactly.
The conclusion nevertheless fails at $-1$.  The map is admissibly
nonmonotone, since $-1\le0$ but $T_n(-1)=1>T_n(0)=0$.

\section*{Why this is the sharp literal obstruction}

The counterexample does not exploit a badly behaved or discontinuous map: it
is an isometry in norm and a 1-Lipschitz lattice operation.  More importantly,
it fixes the entire positive cone.  Consequently no finite or infinite test
family consisting only of nonnegative functions can yield identity
convergence for all continuous sublinear contractions.  Wulbert's linearity
silently supplies oddness, turning information about positive tests into
information about their negatives; plain sublinearity does not.

A viable repair must therefore add sign information or impose asymptotic
oddness.  Indeed, sublinearity always gives
\[
 0\le T_n(h)+T_n(-h).
\]
If both terms converge weakly to $h$ and $-h$, respectively, this nonnegative
defect converges to zero in $L^1$ norm after testing against the constant one.
This eliminates the modulus map on tested directions, but it does not by
itself recover Wulbert's propagation from the test set to all of $L^1$.

\section*{Upgrade and novelty audit}

Six focused routes were checked: the literal modulus construction, arbitrary
positive test families, the sign-symmetric defect, maxima of linear
contractions, scaled/translated corrections, and later-literature searches.
Attempts to preserve both signs with a nonidentity norm-one maximum map run
into Wulbert rigidity and the loss of the $L^1$ norm bound.  No stronger
sign-symmetric proof or counterexample is promoted.

The exact source sentence and combinations of Wulbert, continuous sublinear,
$L^1$, non-monotone, and modulus were searched in the local corpus, arXiv, and
bounded citation neighborhoods through 2026-08-11.  The later
arXiv:2403.03476 generalizes Wulbert within linear operators and confirms the
three hypotheses, but does not treat sublinear maps.  No indexed prior use of
this obstruction for the 2022 question was found; novelty remains subject to
specialist review.

\begin{thebibliography}{9}
\bibitem{GN}
S. G. Gal and C. P. Niculescu,
\emph{Korovkin type theorems for weakly nonlinear and monotone operators},
arXiv:2206.14102, PDF page 12.

\bibitem{W}
D. E. Wulbert,
\emph{Convergence of operators and Korovkin's theorem},
J. Approx. Theory 1 (1968), 381--390, Theorem 9.

\bibitem{KNP}
V. B. Kiran Kumar, M. N. N. Namboodiri, and P. C. Vinaya,
\emph{Korovkin-type approximation for non-positive operators},
arXiv:2403.03476v2, Theorem 1.3.
\end{thebibliography}

\end{document}
